%! Representing a Categorical Variable — Unit 1, Lesson 1.2 — Student Packet Cover Sheet
\documentclass[12pt]{article}
\usepackage{apstats-article}
\usepackage{apstats-boxes}
\usepackage{ltablex}
\keepXColumns

\begin{document}

% Full-bleed navy banner
\begin{tikzpicture}[remember picture, overlay]
  \fill[navy] (current page.north west)
    rectangle ([yshift=-0.9in]current page.north east);
\end{tikzpicture}
\vspace{-0.8in}

\begin{center}
  {\color{white}\textbf{\LARGE Statistics}}\\[4pt]
  {\color{skymid}\large\textbf{Unit 1}}\\[6pt]
  {\color{white}\textbf{\large Lesson 1.2 \quad Representing a Categorical Variable}}
\end{center}

\vspace{0.22in}
\namedateperiod
\vspace{0.18in}

\boxguard[14]
\begin{learningtargetbox}
\small
% GRADUAL RELEASE: the vocabulary is taught in the guided notes, so the targets name it
% plainly. The AP Learning Objectives (UNC-1.A through UNC-1.E) stay in the teacher-facing plan.
\begin{enumerate}[leftmargin=1.5em, itemsep=5pt]
  \item \ldots{} build a \textbf{frequency table} and a \textbf{relative frequency table} for a
        categorical variable, and say what each column has to add up to.
  \item \ldots{} read a \textbf{bar graph} and say exactly which number each bar height is
        showing.
  \item \ldots{} explain why two groups of \textbf{different sizes} can only be compared with
        relative frequencies --- and name the times a count is still the right number.
  \item \ldots{} decide whether a claim like ``most people chose this'' is actually supported
        by the numbers.
\end{enumerate}
\end{learningtargetbox}

\vspace{0.14in}

\boxguard[14]
\begin{tocbox}
\small
\renewcommand{\arraystretch}{1.5}
\begin{tabularx}{\linewidth}{@{} c l X r @{}}
\rowcolor{sky}
\textbf{\#} & \textbf{Component} & \textbf{Description} & \textbf{Score} \\
\midrule
1 & Warm-Up & Percents, and two teams with different numbers of games
                 & \blank{1.2cm} \\
2 & Guided Notes \& Practice & Two tables from one column, what a bar height is, and the
                 comparison that goes wrong when the groups are different sizes --- I show you,
                 we do one together, you do the rest on your own & \blank{1.2cm} \\
% AP Practice is assigned but NOT scored: its score cell prints NA, never a \blank{}.
% Row 4 is the lesson's GRADED practice. Paper homework is the DEFAULT for this course; on the
% occasions DeltaMath is assigned instead, that replaces row 4 and its score cell is struck.
3 & AP Practice & A city recycling audit and an activity-period survey: AP-style multiple
                 choice and one free-response set --- practice, not a grade & \textbf{NA} \\
4 & Homework & A club vote, two hospitals, two versions of one bar graph, and one
                 multiple-choice item --- \textbf{scored, due next class} & \blank{1.2cm} \\
\midrule
  & \multicolumn{2}{r}{\textbf{Total}} & \blank{1.2cm} \\
\end{tabularx}
\end{tocbox}

\vspace{0.14in}

\boxguard[8]
\begin{remindbox}
\small
Today runs in two moves inside the guided notes --- \textbf{I show you}, then \textbf{we do one
together} --- and a third at the end: \textbf{you start the homework here, on your own}, before
you leave. The notes give you the two tables, the bar graph, and the rule for comparing groups of
different sizes; the homework is where you run that rule by yourself.
\textbf{Two questions in this packet have two defensible answers}
depending on which number you use --- when that happens, say which number you used and why.
\end{remindbox}

\end{document}
